\section{Complexity Analysis}
\label{sec:complexity}

\subsection{Per-Component Costs}

\begin{table}[H]
\centering
\small
\caption{Time complexity by component; $n$ = operations in the graph,
$k$ = constraints, $p$ = candidate plans, $d$ = metric dimensions.}
\label{tab:complexity}
\begin{tabular}{@{}lll@{}}
\toprule
Component & Complexity & Notes \\
\midrule
Single constraint (per plan) & $O(n)$ & Safety $O(n^2)$ worst case, $O(n)$ observed \\
Validation (total) & $O(knp)$ & $k = 7$ by default \\
Plan generation & $O(N_{\max} \cdot n)$ & bounded by $N_{\max} = 64$ \\
Cost evaluation & $O(dp)$ & $d = 5$ \\
Dispatch & $O(n)$ & single plan, one-time \\
\midrule
Overall & $O(knp + dp + n)$ & dominated by $O(knp)$ \\
\bottomrule
\end{tabular}
\end{table}

\subsection{Worst Case vs.\ Practice}

The worst case assumes every plan survives all $k$ constraints. With
the default ordering (cheap, high-rejection constraints first),
early-exit reduces effective cost by $2.5$--$3.5\times$: shape and
backend-capability checks reject 30--50\% of candidates, and on average
only 2--3 of 7 constraints execute per plan.

The empirical measurements of Section~\ref{subsec:empirical} bound the
constant factors: for $k = 7$ and $p = 6$ real ResNet-50 plans,
validation plus cost-minimal selection completes in
$1.4$--$8.2\,\mu$s---seven orders of magnitude below a single model
inference on the same machine. The dominant planning cost in practice
is not the algorithm but optional per-candidate \emph{calibration}
(one real execution per plan) when empirically-grounded metrics are
desired.

\subsection{Memory}

$O(p(n + d))$ for candidate plans and metric vectors; with $p \leq 64$
and $d = 5$ this is linear in graph size.
